\documentclass[10pt,a4paper]{amsart}
\usepackage[margin=19mm]{geometry}
\usepackage{amsmath,amssymb,mathtools}
\usepackage[scr=rsfs,cal=euler]{mathalfa}
\usepackage{xfrac}
\usepackage[unicode,hidelinks,bookmarks=false]{hyperref}
\newcommand{\ie}{i.e.\ }
\newcommand{\cf}{cf.\ }
\newcommand{\resp}{respectively\ }
\newcommand{\Subsection}{\subsection}
\providecommand{\nobreakdash}{\nobreak\hbox{-}}
\setlength{\emergencystretch}{2em}
\multlinegap0pt
\usepackage{bm}
\usepackage{bbm}
\usepackage{dsfont}
\usepackage{xspace}
\usepackage{cancel}
\usepackage{mathtools}
\usepackage[shortlabels]{enumitem}
\usepackage{amsthm}
\makeatletter
\@ifundefined{thm}{\newtheorem{thm}{Theorem}}{}
\makeatother
\newcommand{\N}{\mathbb{N}}
\newcommand{\cyc}{\operatorname{cyc}}
\newcommand{\id}{\textit{e}}
\newcommand{\mpair}[1]{\langle\, #1\,\rangle}  % for group generated by S with relations or scalar product
\title[Weak order on groups generated by involutions]{Weak order on groups generated by involutions}
\author{Fabricio Dos Santos, Christophe Hohlweg, Aleksandr Trufanov}
\date{Source: arXiv:2604.18822v1 (CC BY 4.0)}
\begin{document}
\maketitle

\noindent\emph{Source and licence.} Weak order on groups generated by involutions. Version arXiv:2604.18822v1 (CC BY 4.0), distributed under the Creative Commons Attribution 4.0 International licence, as recorded in the arXiv licence metadata for this exact version and shown on the abstract page https://arxiv.org/abs/2604.18822v1. Full rights notice: licenses/arxiv-2604.18822-CC-BY-4.0.txt.\par\smallskip

\noindent\textit{Editorial scope.} Counted unit: the Thm whose source label is \texttt{thm:rk3-reco} in the article (source0.tex lines 1330--1343, source label \texttt{thm:rk3-reco}), delivered together with the article's own complete original proof of it (source0.tex lines 1345--1372, one \texttt{proof} environment of 28 lines). No mathematics is written, repaired or re-derived: statement and proof are the article's own text line for line. Required context retained uncounted: condition 1 - recognizable (lines 1277--1287); condition 2 - recognizable (lines 1277--1287); condition 3 - recognizable (lines 1277--1287); condition 4 - recognizable (lines 1277--1287). Every definition, display and label of those lines is retained unchanged. Omitted from this package and not redistributed: 1--1276, 1288--1329, 1344--1344, 1373--2309. Nothing inside a retained excerpt is omitted or altered: every retained body is delivered exactly as the article's lines are written, so no omission receipt of any kind accompanies this package. Citations: the retained excerpts carry no citation command at all, so no bibliography entry of the article is transcribed and no reference is redistributed. Reference adaptation: none was needed and none was made; every reference inside the delivered unit resolves inside it, so no unresolved reference or citation is produced. Printed slips kept literal: no printed word, symbol, hypothesis or number of the counted statement or of its proof was altered. Layout and class: the article's own class is \texttt{amsart} with packages including inputenc, hyperref, amsmath, amssymb, amsthm, mathtools, xspace, xcolor, tikz, pgfplots, array, graphicx, caption, booktabs; the wrapper is the lane's pinned \texttt{amsart} editorial wrapper at 10pt on A4, which restates the theorem-like environments the retained text needs and adds the article's own macro definitions that the retained text needs (\texttt{\textbackslash N}, \texttt{\textbackslash cyc}, \texttt{\textbackslash id}, \texttt{\textbackslash mpair}), with extra wrapper packages (bm, bbm, dsfont, xspace, cancel, mathtools, enumitem). The delivered numbering is the unit's own. Assets: no retained span contains a figure environment, a graphics-inclusion command, a PDF-page inclusion, a table or a TikZ picture; no image file is copied into this package, the package declares no asset dependency and no image file is redistributed with it. Licence: version arXiv:2604.18822v1 of this article is distributed under the Creative Commons Attribution 4.0 International licence, as recorded in the arXiv licence metadata for this exact version and shown on the abstract page https://arxiv.org/abs/2604.18822v1. A full rights notice is delivered at licenses/arxiv-2604.18822-CC-BY-4.0.txt. Characters: every retained excerpt is pure ASCII, so the delivered engine, XeLaTeX with its default Latin Modern text font, reports no missing character.
\begin{enumerate}[series=2-recognizable]    
    \item \label{condition 1 - recognizable} $[w]$ is a reduced cyclic word for every $w \in R_0$;

    \item \label{condition 2 - recognizable} Every $w \in R_0$ has even length at least $4$;

	\item
	\label{condition 3 - recognizable}
	Let $w\in R_0$ and $ab$ a $2$-factor of $[w]$. If $ab$ is also $2$-factor of $[v]$ or $[v^{-1}]$ for some $v\in R_0$, then $w=v$.
   
    \item \label{condition 4 - recognizable} Let $w=s_1\cdots s_k\in R_0$ and $ab$ be a $2$-factor of $s_kws_1$. If $s\in S$ precedes (resp. follows) $ab$ in $s_{k-1} s_kws_1 s_2$, then for any occurrence of the factor $ab$ in $s_kws_1$,  $s$  also precedes (resp. follows) that occurence in $s_{k-1} s_kws_1 s_2$.
\end{enumerate}

\begin{thm}
\label{thm:rk3-reco}
    Let $(W,S)$ be an involution system of rank 3. Then $(W,S)$ is a $2$-recognizable involution system if and only if it admits one of the following presentations (where $\infty$ means there is no relation): 
    \begin{enumerate}[(i)]
        \item $\langle a,b,c \, | \, a^2=b^2=c^2 = (abc)^{2m}=\id\rangle$, $m \in \N_{\geq 1}$; \label{type1}
        
        \item $\langle a,b,c \, | \, a^2=b^2=c^2 = (abacbc)^{m}=\id \rangle$, $m \in \N_{\geq 1}$; \label{type2}
        
        \item $\langle a,b,c \, | \, a^2=b^2=c^2 = (abac)^m=(bc)^n=\id\rangle$, $m \in \N_{\geq 1}$, $n \in \N_{\geq 2} \cup \{\infty\}$; \label{type3}

        \item $\langle a,b,c \, | \, a^2=b^2=c^2 = (ab)^m=(bc)^n=(ac)^l=\id\rangle$, $m,n,l \in \N_{\geq 2} \cup \{\infty\}$ (the case of Coxeter systems of rank $3$). \label{type4}
    
    \end{enumerate}    
\end{thm}

\begin{proof}  It is readily seen that the presentations given above are $2$-recognizable. So we only need to classify all 2-recognizable presentations $W_R = \mpair{S \mid \{s^2 \mid s \in S\} \cup R_0}$ where $|S|=3$. We do this up to renaming letters or replacing an element $w \in R_0$ by $w' \in \cyc(w)$ since these operations do not change the group up to isomorphism.
    
    
    \smallskip
    Let $S = \{a,b,c\}$. Since there are only three letters in $S$, we must have that $|R_0| \leq 3$ by (\ref{condition 3 - recognizable}). If $R_0=\emptyset$, we get presentation \ref{type4} with $m=n=l=\infty$.

    \smallskip
    \noindent {\bf Case $|R_0| = 1$}.  Let $w = s_1 \cdots s_k \in R_0$ be nonempty. Without loss of generality we can assume $w$ starts with $ab$.
    
    Consider first the case when $xyx$ is not a subword of $w$, where $x,y \in S$. This forces the third letter of $w$ to be $c$, the fourth to be $a$, the fifth to be $b$, and so on. Since $w$ is cyclically reduced by (\ref{condition 1 - recognizable}) it cannot end with $a$. Furthermore, $w$ cannot end in $b$, otherwise $s_kws_1 = bwa$ would have an occurrence of $ab$ followed by $c$ and another followed by $a$, contradicting (\ref{condition 4 - recognizable}). Hence, since $w$ has even length at least 4 by (\ref{condition 2 - recognizable}), we must have that $w=(abc)^{2m}$ for some $m \in \N_{\geq 1}$ and we get presentation \ref{type1}. 
    
    Now, consider the case when $xyx$ is a subword of $w$ for some $x,y \in S$. Without loss of generality, assume $w$ starts with $aba$. 
    
    If the fourth letter of $w$ is $b$, then $w$ cannot contain the letter $c$. Otherwise, we could either find two occurrences of $ab$ in $s_kws_1$ where one is followed by $a$ and the other by $c$, or two occurrences of $ba$ in $s_kws_1$ where one is followed by $b$ and the other by $c$; contradicting (\ref{condition 4 - recognizable}). Since $w$ is cyclically reduced by (\ref{condition 1 - recognizable}) it cannot end with $a$, so in this case we get presentation \ref{type4} with $n=l=\infty$. 
    
    On the other hand, if the fourth letter of $w$ is $c$, then we have three cases. The first case is when $w$ has length 4, so $w=abac$ and we obtain the presentation \ref{type3} with $n=\infty$. 
    
    The second case is $|w|\geq 5$ and the fifth letter is $a$. We then claim that $w = (abac)^m$ for some $m\in \N_{\geq 1}$. Suppose not, that is $w = (abac)^mu$ for some $u \neq abac$ nonempty (note that if $m=1$ then $u$ starts with $a$). Let $s$ be the first letter of $u$ differing from $abac$, that is $u=u_1su_2$ where $u_1$ is a prefix of $abac$ but $abac \neq u_1su_2'$ for any word $u_2'$ (note that $s$ is possibly empty if $|u|<4$). Let $r_1,r_2$ be the two letters preceding $s$ in $w$. Then $r_1r_2$ is a 2-factor of $s_kws_1$ that has two occurrences in $s_kws_1$ followed by different letters, contradicting (\ref{condition 4 - recognizable}). Hence, we get presentation \ref{type3} with $n=\infty$. 
    
    Finally, the third case is $|w|\geq 5$ and the fifth letter is $b$. First, note that $w$ cannot be $abacb$ since $s_kws_1 = babacba$ has an occurrence of $ba$ followed by $b$ and another followed by $c$ and this would contradict (\ref{condition 4 - recognizable}). Hence, $|w|\geq 6$ and the sixth letter has to be $c$ since $w$ is cyclically reduced by (\ref{condition 1 - recognizable}). Then, a similar argument to the one in the previous paragraph gives us that $w = (abacbc)^m$ for some $m\in \N_{\geq 1}$, so we get presentation \ref{type2}.


    \smallskip
    \noindent {\bf Case $|R_0| = 2$}. Let $w_1,w_2 \in R_0$ be nonempty. If $w_1$ and $w_2$ only contain two different letters, then we get the presentation \ref{type4} with one of $m,n,l$ being infinity. Otherwise, assume that $w_1$ contains three different letters. From the case $|R_0|= 1$ we know that $w_1$ must be of the form $(abc)^{2m}$, $(abacbc)^m$ or $(abac)^m$ for some $m \in  \N_{\geq 1}$. But by (\ref{condition 3 - recognizable}), $w_1$ cannot be $(abc)^{2m}$ or  $(abacbc)^m$ since $xy$ or $yx$ is a 2-factor of $[(abc)^{2m}]$ and $[(abacbc)^m]$ for any $x,y \in S$. Hence, we must have $w_1 = (abac)^m$. Then, the only possibility for $w_2$ is $w_2 = (bc)^n$ where $n \in  \N_{\geq 2}$. Hence, we get presentation \ref{type3}.

    \smallskip
    \noindent {\bf Case $|R_0| = 3$}. Let $w_1,w_2,w_3 \in R_0$ be nonempty. If all of them only contain two different letters, then we get the presentation \ref{type4}. Otherwise, assume that $w_1$ contains three different letters. From the case $|R_0| = 2$, we must have that $w_1 = (abac)^m$ for some $m \in  \N_{\geq 1}$ and $w_2 = (bc)^n$ where $n \in  \N_{\geq 2}$. But if $w_3$ is nonempty, then we get a contradiction to (\ref{condition 3 - recognizable}) since $xy$ or $yx$ is a 2-factor of $[(abac)^m]$ or $[(bc)^n]$ for any $x,y \in S$. Hence, no word in $R_0$ can have three different letters, so the only possibility is presentation \ref{type4}.
  \end{proof}

\end{document}
